\chapter{Arithmetic Bézout and intersection estimates}
\label{chap:arithbezout}

% archon:covers MR4276287UniformityInMordellLangForCurves/ArithBezout.lean

This chapter collects the arithmetic and analytic intersection-theoretic inputs
to the auxiliary height inequality \cref{prop:aux_ht_ineq}. The centrepiece is
the arithmetic Bézout theorem of Autissier (\cref{thm:arith_bezout}), used as a
black box. The remaining lemmas establish the bihomogeneous degree bound for the
graph of \([N]\) (\cref{lem:bihomog_degree}) and the two intersection-number
estimates (\cref{lem:siu_intersection_lower,lem:siu_intersection_upper}) whose
comparison, via Siu's bigness criterion, yields the height inequality.

\section{Arithmetic Bézout theorem}

\begin{theorem}[Arithmetic Bézout theorem]
  \label{thm:arith_bezout}
  \lean{MR4276287UniformityInMordellLangForCurves.ArithBezout.arithBezout}
  \uses{def:ambient_setup}
  % SOURCE: [Théorème 3]{HauteursAlt3} (verbatim text not yet retrieved)
  \textit{Source: Autissier, \cite[Théorème~3]{HauteursAlt3}.}
  Let \(Z_1, Z_2\) be algebraic cycles on \(\mathbb{P}^n_{\mathbb{Z}}\) of
  complementary codimension. Their arithmetic intersection number satisfies
  \[
    \hat{h}(Z_1 \cdot Z_2)
    \le \deg(Z_1)\,\hat{h}(Z_2)
       + \deg(Z_2)\,\hat{h}(Z_1)
       + c\,\deg(Z_1)\,\deg(Z_2),
  \]
  where \(c > 0\) is an absolute constant and \(\hat{h}\) denotes the
  normalized height of a cycle.
\end{theorem}

\begin{proof}
  \uses{def:ambient_setup}
  This is \cite[Théorème~3]{HauteursAlt3} (Autissier). The proof is an arithmetic
  intersection theory computation in the Arakelov geometry of \(\mathbb{P}^n_{\mathbb{Z}}\),
  using the arithmetic Chern class formalism. The bound is a consequence of the
  Cauchy--Schwarz inequality for the arithmetic intersection pairing after normalizing
  the Green functions at the archimedean places. We axiomatize this theorem as a
  black box and do not reprove it in this project.
\end{proof}

\section{Bihomogeneous degree of the graph of \([N]\)}

\begin{lemma}[Bihomogeneous degree bound for \([N]\)]
  \label{lem:bihomog_degree}
  \lean{MR4276287UniformityInMordellLangForCurves.ArithBezout.bihomogDegree}
  \uses{lem:graph_construction}
  % SOURCE: Section 4, proof of Proposition~4.2 of \cite{MR4276287-mordell-lang-curves} (read from references/MR4276287-mordell-lang-curves-source/HtIneq.tex)
  % SOURCE QUOTE: "If for all $i$ the polynomials $f^{(l)}_i$ are bihomogeneous of degree $(D_{l},D'_{l})$, then all $f^{(l+1)}_i$ are bihomogeneous of degree $(4D_{l},D'+4D'_{l})$. Recall that $(D_1,D'_1)=(4,D')$, thus the recurrence relations $D_{l+1}=4D_l  \quad\text{and}\quad D'_{l+1} = D' + 4D'_{l}$ imply $D_l = 4^l \quad\text{and}\quad   D'_l = \frac{4^l-1}{3}D'\le 4^lD'$ for all $l\ge 1$."
  \textit{Source: Section 4, proof of Proposition~4.2 of \cite{MR4276287-mordell-lang-curves} (read from references/MR4276287-mordell-lang-curves-source/HtIneq.tex).}
  Let \(N = 2^l\). The graph of multiplication-by-\(N\) on \(\mathcal{A}\),
  viewed as a closed subvariety of
  \(\mathbb{P}^n \times \mathbb{P}^n \times \mathbb{P}^m\) via the graph
  closure \(\overline{X_N}\) from \cref{lem:graph_construction}, is cut out
  by \(n\) equations that are trihomogeneous of tridegree \((D_l, 1, D'_l)\)
  where
  \[
    D_l = 4^l = N^2, \qquad D'_l = \tfrac{4^l - 1}{3}\,D' \le D' N^2,
  \]
  for an absolute constant \(D'\). In particular, the degree in each of the
  \(\mathbf{X}\)- and \(\mathbf{S}\)-blocks grows at most as \(N^2\).
\end{lemma}

% SOURCE QUOTE PROOF: "For each integer $l\ge 1$ we require polynomials $f^{(l)}_0,\ldots,f^{(l)}_n$ to describe multiplication-by-$2^l$, \textit{cf.} \cite[$\mathsection$9]{GaoHab}. In order to obtain information on the degree with respect to $\mathbf{S}$ we construct them by iterating the $f^{(1)}_0=f_0,\ldots,f^{(1)}_n=f_n$. For all $i\in \{0,\ldots,n\}$ we set $f^{(l+1)}_i(\mathbf{X},\mathbf{S}) = f_i\left(\bigl(f^{(l)}_0(\mathbf{X},\mathbf{S}),\ldots,f^{(l)}_n(\mathbf{X},\mathbf{S})\bigr),\mathbf{S}\right)$ for all $i$; it is bihomogeneous in $\mathbf{X}$ and $\mathbf{S}$."
\begin{proof}
  \uses{lem:graph_construction}
  Multiplication-by-\(2\) on the projectively normal abelian variety
  \(A \subset \mathbb{P}^n_{k(S)}\) satisfies \([2]^* L \cong L^{\otimes 4}\) for the
  symmetric ample line bundle \(L\); by a result of Serre
  \cite[Corollaire~2, Appendix~II]{Ast6970}, \([2]\) is therefore represented by
  bihomogeneous polynomials \(f_0, \ldots, f_n\) in the projective coordinates
  \(\mathbf{X}\) of \(\mathbb{P}^n\) and the base coordinates \(\mathbf{S}\) of
  \(\mathbb{P}^m\), of bidegree \((4, D')\) for some constant \(D'\).
  Define \(f^{(1)}_i = f_i\) and iterate:
  \[
    f^{(l+1)}_i(\mathbf{X}, \mathbf{S})
    = f_i\bigl((f^{(l)}_0(\mathbf{X}, \mathbf{S}),
               \ldots,
               f^{(l)}_n(\mathbf{X}, \mathbf{S})),\, \mathbf{S}\bigr).
  \]
  If \(f^{(l)}_i\) has bidegree \((D_l, D'_l)\), then \(f^{(l+1)}_i\) has
  bidegree \((4D_l,\, D' + 4D'_l)\). Starting from \((D_1, D'_1) = (4, D')\),
  the recurrence gives \(D_l = 4^l\) and \(D'_l = \tfrac{4^l - 1}{3}D' \le D' \cdot 4^l\).
  The graph of \([N] = [2^l]\) is then cut out inside \(\overline{X} \times \mathbb{P}^n\)
  by the \(n\) trihomogeneous polynomials
  \(g_i = Y_i f^{(l)}_0(\mathbf{X}, \mathbf{S}) - Y_0 f^{(l)}_i(\mathbf{X}, \mathbf{S})\)
  of tridegree \((D_l, 1, D'_l) = (N^2, 1, D'_l)\).
\end{proof}

\section{Lower bound on \((\mathcal{F}^{\cdot d})\)}

\begin{lemma}[Lower bound on the self-intersection of \(\mathcal{F}\)]
  \label{lem:siu_intersection_lower}
  \lean{MR4276287UniformityInMordellLangForCurves.ArithBezout.siuIntersectionLower}
  \uses{prop:betti_form,lem:graph_construction,def:isNef}
  % SOURCE: Proposition 4.1, Section 4 of \cite{MR4276287-mordell-lang-curves} (read from references/MR4276287-mordell-lang-curves-source/HtIneq.tex)
  % SOURCE QUOTE: "Suppose $\an{X}$ contains a smooth point at which $\omega|_{\sman{X}}^{\wedge d} >0$. Then there exists a constant $\kappa > 0$, independent of $N$, such that $(\cF^{\cdot d}) \ge \kappa N^{2d}$ for all $N\in\IN$."
  \textit{Source: Proposition 4.1, Section 4 of \cite{MR4276287-mordell-lang-curves} (read from references/MR4276287-mordell-lang-curves-source/HtIneq.tex).}
  Suppose \(X^{\mathrm{an}}\) contains a smooth point at which
  \(\omega|_{X^{\mathrm{sm,an}}}^{\wedge d} > 0\),
  where \(\omega\) is the Betti form from \cref{prop:betti_form}.
  Then there exists a constant \(\kappa > 0\), independent of \(N\), such that
  \[
    (\mathcal{F}^{\cdot d}) \ge \kappa N^{2d}
  \]
  for all \(N \in \mathbb{N}\), where
  \(\mathcal{F} = \mathcal{O}(0,1,1)|_{\overline{X_N}}\) is the line bundle on the
  graph closure \(\overline{X_N} \subset \mathbb{P}^n \times \mathbb{P}^n \times \mathbb{P}^m\)
  introduced in \cref{lem:graph_construction}.
\end{lemma}

% SOURCE QUOTE PROOF: "We fix a point $P_0 \in \sman{X}$ at which $\omega|_{\sman{X}}^{\wedge d}$ is positive and let $s_0=\pi(P_0),\Delta,\vartheta,\theta,$ and $K$ be as in $\mathsection$\ref{SubsectionSetUpHtIneqBetti}, see Remark~\ref{rmk:nondeg}. In particular $\vartheta(P_0) = \theta\circ \pi(P_0) = 1$. We extend $\vartheta$ to a smooth function on $\an{(\IP^m_\IC)}$ by setting it $0$ outside of the compact set $K\subset \an{S}$. This extends $\theta = \vartheta\circ\pi$ to all of $\an{(\IP^n_\IC\times\IP^m_\IC)}$."
\begin{proof}
  \uses{prop:betti_form,lem:graph_construction}
  Fix a smooth point \(P_0 \in X^{\mathrm{sm,an}}\) at which
  \(\omega|_{X^{\mathrm{sm,an}}}^{\wedge d} > 0\). Let \(\alpha\) be the pull-back
  of the Fubini--Study form on \(\mathbb{P}^{(n+1)(m+1)-1}_\mathbb{C}\) under the
  analytification of the Segre morphism
  \(\mathbb{P}^n_\mathbb{C} \times \mathbb{P}^m_\mathbb{C} \to \mathbb{P}^{(n+1)(m+1)-1}_\mathbb{C}\),
  restricted to \(\overline{\mathcal{A}}^{\mathrm{an}}\); it represents
  \(\mathcal{O}(1,1)|_{\overline{\mathcal{A}}^{\mathrm{an}}}\) and is strictly positive on
  \(\mathcal{A}^{\mathrm{an}}\).
  Since the Betti form \(\omega\) is semi-positive by \cref{prop:betti_form}(i),
  there exist a compact neighbourhood \(K \subset \Delta \subset \pi(P_0)\) and
  a constant \(C > 0\) such that \(C\alpha|_{\mathcal{A}_\Delta} - \omega|_{\mathcal{A}_\Delta} \ge 0\).
  Let \(\theta\) be a smooth non-negative cutoff supported on \(\pi^{-1}(K)\) with
  \(\theta(P_0) = 1\). Using the Betti-form scaling \([N]^*\omega = N^2\omega\)
  from \cref{prop:betti_form} and the fact that \(\theta = \vartheta \circ \pi\) is
  pulled back from the base (so \([N]^*\theta = \theta\)), one obtains
  \[
    \beta := C [N]^*(\theta\alpha) - N^2\theta\omega \ge 0
    \quad \text{on } \mathcal{A}^{\mathrm{an}}.
  \]
  Writing \(\gamma = C [N]^*(\theta\alpha)\) and \(\delta = N^2\theta\omega\), both
  semi-positive with compact support, a wedge-positivity argument
  \cite[Proposition~III.1.11]{Demailly} gives
  \[
    \int_{X^{\mathrm{sm,an}}} \gamma^{\wedge d}
    \ge \int_{X^{\mathrm{sm,an}}} \delta^{\wedge d}
    = \kappa' N^{2d}, \qquad
    \kappa' = \int_{X^{\mathrm{sm,an}}} (\theta\omega)^{\wedge d} > 0,
  \]
  the strict positivity \(\kappa' > 0\) following from the choice of \(P_0\) and
  \(\theta(P_0) = 1\). Via the graph construction \cref{lem:graph_construction}
  (\(\overline{X_N} \xleftarrow{\rho_1} X_N \xrightarrow{\rho_2} \mathcal{A}\)),
  a change-of-coordinates and the identification
  \(\rho_2^*\mathcal{O}(1,1) = \mathcal{O}(0,1,1)|_{\overline{X_N}} = \mathcal{F}\)
  via the projection formula yield
  \[
    \int_{X^{\mathrm{sm,an}}} [N]^*(\theta\alpha)^{\wedge d} \le (\mathcal{F}^{\cdot d}).
  \]
  Combining, \((\mathcal{F}^{\cdot d}) \ge (\kappa'/C^d)\,N^{2d}\), and
  \(\kappa = \kappa'/C^d > 0\).
\end{proof}

\section{Upper bound on \((\mathcal{M} \cdot \mathcal{F}^{\cdot(d-1)})\)}

\begin{lemma}[Upper bound on the mixed intersection]
  \label{lem:siu_intersection_upper}
  \lean{MR4276287UniformityInMordellLangForCurves.ArithBezout.siuIntersectionUpper}
  \uses{lem:bihomog_degree,lem:graph_construction,def:isNef}
  % SOURCE: Proposition 4.2, Section 4 of \cite{MR4276287-mordell-lang-curves} (read from references/MR4276287-mordell-lang-curves-source/HtIneq.tex)
  % SOURCE QUOTE: "Assume $d \ge 1$. There exists a constant $c>0$ depending on the data introduced above with the following property. Say $N\ge 1$ is a power of $2$, then $(\cM\cdot \cF^{\cdot(d-1)}) \le c N^{2(d-1)}.$"
  \textit{Source: Proposition 4.2, Section 4 of \cite{MR4276287-mordell-lang-curves} (read from references/MR4276287-mordell-lang-curves-source/HtIneq.tex).}
  Assume \(d \ge 1\). There exists a constant \(c > 0\) depending only on \(X\) and
  the ambient data, with the following property. For all \(N \ge 1\) a power of \(2\),
  \[
    (\mathcal{M} \cdot \mathcal{F}^{\cdot(d-1)}) \le c\, N^{2(d-1)},
  \]
  where \(\mathcal{M} = \mathcal{O}(0,0,1)|_{\overline{X_N}}\) and
  \(\mathcal{F} = \mathcal{O}(0,1,1)|_{\overline{X_N}}\) are the line bundles
  on the graph closure \(\overline{X_N}\) from \cref{lem:graph_construction}.
\end{lemma}

% SOURCE QUOTE PROOF: "Recall that $[2] \colon A \rightarrow A$ is the multiplication-by-$2$ morphism on $A$. For the symmetric and ample line bundle $L$ on $A$, we have $[2]^*L \cong L^{\otimes 4}$. Recall that $A$ is projectively normal in $\IP_{k(S)}^n$. By a result of Serre, \cite[Corollaire~2, Appendix~II]{Ast6970}, the morphism $[2]$ is represented by homogeneous polynomials $f_0,\ldots,f_n$ in the $n+1$ projective coordinates of $\IP^n$ of degree $4$, with coefficients in $k(S)$ and with no common zeros in $A$."
\begin{proof}
  \uses{lem:bihomog_degree,lem:graph_construction}
  By \cref{lem:bihomog_degree} with \(N = 2^l\), the graph \(\overline{X_N}\) is an
  irreducible component of the intersection of \(\overline{X} \times \mathbb{P}^n\)
  with the hypersurfaces \(V(g_1), \ldots, V(g_n)\) of tridegree \((N^2, 1, D'_l)\)
  in \(\mathbb{P}^n \times \mathbb{P}^n \times \mathbb{P}^m\), with \(D'_l \le D' N^2\).
  The line bundles \(\mathcal{O}(0,0,1)\) and \(\mathcal{O}(0,1,1)\) are nef on
  \(\mathbb{P}^n \times \mathbb{P}^n \times \mathbb{P}^m\) (any effective divisor on a
  product of projective spaces is nef). Fulton's positivity result
  \cite[Corollary~12.2(a)]{Fulton} gives
  \begin{align*}
    (\mathcal{O}(0,0,1)\,\mathcal{O}(0,1,1)^{\cdot(d-1)}[\overline{X_N}])
    &\le \bigl(\mathcal{O}(0,0,1)\,\mathcal{O}(0,1,1)^{\cdot(d-1)}\,
         \mathcal{O}(N^2,1,D'_l)^{\cdot n}[\overline{X} \times \mathbb{P}^n]\bigr).
  \end{align*}
  Expanding by multilinearity and writing the class \([\overline{X} \times \mathbb{P}^n]\)
  in terms of the hyperplane pullbacks \(H_1, H_2\) of the \(\mathbb{P}^n_\mathbf{X}\)-
  and \(\mathbb{P}^m_\mathbf{S}\)-factors, only terms with total degrees
  \((i, j, p) = (n, n, m)\) in the three hyperplane pullbacks contribute.
  Since \(D_l = N^2\) and \(D'_l \le D' N^2\), the sum of contributing terms
  is bounded by
  \[
    (N^2 D')^{d-1} \cdot C_X
  \]
  for a constant \(C_X\) depending only on \(n\), \(m\), \(d\), and the cycle
  class of \(\overline{X}\). Recalling \(\mathcal{M} = \mathcal{O}(0,0,1)|_{\overline{X_N}}\)
  and \(\mathcal{F} = \mathcal{O}(0,1,1)|_{\overline{X_N}}\), we obtain
  \((\mathcal{M} \cdot \mathcal{F}^{\cdot(d-1)}) \le c N^{2(d-1)}\) with
  \(c = {D'}^{d-1} C_X\).
\end{proof}
